% !TEX root = ../main.tex
% ---------------------------------------------------------------------
% Apêndice B — registro das 31 decisões metodológicas.
%
% ATENÇÃO: este arquivo é MANTIDO À MÃO, ao contrário do 07_apendice.tex, que
% é gerado por apendice/gerar_apendice.py a partir dos CSVs dos pipelines. O
% conteúdo aqui é editorial (o que se decidiu e por quê), não saída de
% pipeline, e por isso não tem gerador.
%
% A fonte espelhada é o bloco #p4-decisoes da Visualizacao/index.html. Quando
% uma decisão nova entrar, ela entra nos dois lugares. O verificar.py confere
% que todo DNN citado no corpo tem entrada aqui; ele NÃO confere o contrário
% (dez das trinta e uma não são citadas em passagem nenhuma, e isso é o
% esperado num registro completo) nem confere o texto de uma contra a outra,
% que é trabalho humano.
%
% ⚠ Número que entra aqui tem de existir no corpo do texto. O rascunho deste
% arquivo trouxe três números do registro da viz que o texto impresso não
% carrega, e um apêndice que introduz número não conferível é pior que a
% citação sem destino que ele veio resolver.
% ---------------------------------------------------------------------
\chapter{Registro das decisões metodológicas}
\label{ap:decisoes}

Cada escolha de desenho que alterou algum resultado já obtido, ou que fixou
uma regra à qual as análises seguintes ficaram submetidas, recebeu um número
no curso da pesquisa, e são trinta e uma ao todo. Onde uma delas incide sobre
uma passagem do texto, a passagem a cita pelo número e não repete ali o que
ela decidiu, e é por isso que este apêndice existe. Vinte e uma são citadas assim;
as dez restantes valem para o trabalho todo sem que nenhuma passagem em
particular dependa delas, e estão aqui porque um registro com buracos deixa de
ser registro.

A numeração segue a ordem em que as decisões foram tomadas, e não a ordem de
importância, o que explica por que uma entrada de uma linha convive aqui com
outra que ocupa um parágrafo inteiro. Ela também não é reaproveitada, de modo
que uma decisão superada permanece no registro com indicação de qual a
substituiu, como acontece com a D18. Manter a decisão superada tem função,
porque é ela que informa o que teria sido publicado se a conferência posterior
não tivesse existido.

Sete episódios de correção alteraram resultados já obtidos e recebem
tratamento próprio na Seção~\ref{sec:met-decisoes}, onde constam o defeito, a
correção e a regra geral que ficou no lugar. Episódio não é sinônimo de
decisão, porque um deles fixou quatro decisões de uma vez, outro fixou duas e
um terceiro não gerou decisão numerada nenhuma, e é por isso que sete
episódios cobrem dez das entradas abaixo. As entradas correspondentes aqui são
resumidas e remetem àquela seção. Também não se confundem com os sete
resultados superados do Quadro~\ref{quadro:superados}, que são outra lista e
respondem a outra pergunta, a de o que teria sido publicado sem o teste que o
superou. O mesmo registro, em versão navegável e acompanhado das séries que
cada decisão moveu, está na visualização interativa (\url{\urlviz}).

\section{Classes e medição}

\decisao{D1}{Seis classes nos estoques e sete nos fluxos, com o Mosaico de
Usos como grupo próprio}{O agrupamento reduz o ruído das subclasses e é o
mesmo em mapas, matrizes de transição e taxas. O Mosaico ganhou identidade
própria nos fluxos porque é ele que absorve a conversão no fim da série, e
continua fora dos mapas anuais de cobertura, de modo que a regra é que o fluxo
o pinta e o estoque não.}

\decisao{D2}{As taxas vêm do MapBiomas diretamente}{A fonte é a série
municipal original, e não o painel derivado, o que impede que uma etapa
intermediária de agregação redefina a medida sem que isso apareça.}

\decisao{D10}{A idade da pastagem é calculada localmente, em Python}{O cálculo
percorre as quarenta bandas anuais, e o encadeamento de mais de trinta
operações que ele exige estoura o limite de complexidade do Google Earth
Engine.}

\decisao{D13}{``Terra convertível'' é um \emph{proxy} do MapBiomas, lido como
teto}{Não existe medida direta da oferta de Cerrado passível de conversão, e o
que se pode construir é um limite superior. Ele é calculado em três definições
alternativas, para que a conclusão não dependa de qual delas se adote.}

\decisao{D17}{A proteção legal entra como malha vetorial, e não como classe de
cobertura}{As unidades de conservação vêm do Cadastro Nacional e as terras
indígenas da Funai, e ambas são cruzadas por sobreposição com as AMCs. Foi
assim que se mediu quanto do Cerrado convertível remanescente está fora de
proteção integral, que é a única categoria a vedar a conversão de modo geral,
entre 94,3\% e 97\% conforme a conta seja refeita pixel a pixel ou feita sobre
a malha de unidades (Seção~\ref{sec:res-perna4}). Estar fora dela não torna o
restante livre, porque Reserva Legal e Área de Preservação Permanente
continuam incidindo sobre ele.}

\decisao{D18}{O custo de carbono é reportado em três cenários de densidade, e
nunca em número único}{As densidades vinham compiladas da literatura sob o
nível 1 do IPCC, e o resultado era reportado como faixa, e não como ponto, o
que reconhecia a incerteza do parâmetro sem resolver a procedência dele. A
decisão fica registrada, e não removida, porque foi superada pela D30 e é o
que se teria publicado antes da conferência que a substituiu
(Seção~\ref{sec:met-d30}).}

\decisao{D30}{Havendo inventário oficial para o parâmetro, é dele que o
parâmetro vem}{Os estoques de carbono passaram a ser os do Quarto Inventário
Nacional, que os publica por fitofisionomia e com valor próprio para Goiás.
Da troca saiu também a regra de que, quando uma conclusão depende de uma razão
entre parâmetros, é a razão que precisa entrar na análise de sensibilidade (Seção~\ref{sec:met-d30}).}

\decisao{D31}{Estoque removido não é emissão líquida, e as duas grandezas são
calculadas e reportadas com nomes distintos}{Atribuir a cada hectare
convertido a densidade da formação que ali estava mede o que saiu da cobertura
de origem, e o método de diferença de estoque pede também o que o uso entrante
repõe (Seção~\ref{sec:met-d31}).}

\section{Tempo, taxas e janelas}

\decisao{D3}{A inclinação da tendência é estimada em janela móvel de cinco
anos}{São mantidas duas versões, uma retroativa, que é a usada em inferência,
e uma centrada, que serve apenas à visualização e nunca a um teste.}

\decisao{D4}{O erro-padrão dessas inclinações é o de Newey-West, com duas
defasagens}{Janelas móveis sobrepostas produzem autocorrelação de resíduos por
construção, e um erro-padrão que a ignore é otimista por desenho.}

\decisao{D5}{A aceleração é a diferença entre inclinações consecutivas}{É a
métrica mais frágil do conjunto, porque acumula o ruído de duas estimativas,
e por isso só entram nos gráficos os pontos que superam dois desvios-padrão.}

\decisao{D12}{As janelas temporais são explícitas, e toda análise por período
roda em quatro resoluções}{As quatro são a série contínua, os três atos, uma
grade regular de cinco anos e as décadas, e o que se reporta é a concordância
entre elas. É também o princípio que impede uma análise de reimportar as
fronteiras da periodização que ela deveria testar (Seção~\ref{sec:met-robustez}).}

\section{Espaço e unidades de análise}

\decisao{D6}{A malha regional de base é a das cinco mesorregiões do IBGE de
2017}{Não a das 22 regiões geográficas imediatas, porque é sobre as
mesorregiões que o eixo Sul--Norte se lê com unidades grandes o bastante para
comportar a comparação entre regimes, e porque a análise fina já corre sobre
as 166 AMCs. Parte dos resultados agrega as cinco em três blocos latitudinais,
o Sul, o Centro e o Norte, onde o número de unidades por bloco precisa ser
maior, e o texto declara qual das duas malhas está em uso a cada vez
(Seção~\ref{sec:met-area}).}

\decisao{D11}{As séries longitudinais são agregadas em 166 Áreas Mínimas
Comparáveis de território constante}{Um quarto dos municípios goianos de hoje
foi criado depois de 1985, e a série municipal registraria como queda o que é
apenas desmembramento de território (Seção~\ref{sec:met-amc}).}

\section{Inferência e identificação}

\decisao{D7}{Nenhuma associação é lida em nível, e todas passam para primeiras
diferenças}{Duas séries com tendência comum correlacionam-se sem que exista
relação entre elas, e a primeira diferença é o que retira essa tendência antes
de qualquer leitura (Seção~\ref{sec:met-painel}).}

\decisao{D8}{O painel tem efeitos fixos de unidade e de ano, com erro-padrão
agrupado}{Os efeitos de unidade absorvem a heterogeneidade local que não varia
no tempo, e os de ano absorvem os choques comuns a todas as unidades no mesmo
ano (Seção~\ref{sec:met-instrumental}).}

\decisao{D9}{O desenho de diferenças em diferenças entre Goiás, Mato Grosso e
Tocantins é lido como sensibilidade, e nunca como causa}{Os quatro marcos
testados são federais e incidem no mesmo dia sobre o estado tratado e sobre os
controles, de modo que não existe grupo não tratado e o coeficiente mede
exposição diferencial a um choque comum. Uma ressalva posterior acrescenta que
tratado e controle nem sequer partem do mesmo piso de Reserva Legal, porque o
cerrado tocantinense está na Amazônia Legal e Goiás não, o que é mais uma
razão para não ler o contraste entre os dois como efeito de lei. Nenhuma
conclusão deste trabalho se apoia neste desenho
(Seção~\ref{sec:met-camadas}).}

\decisao{D14}{Em recorte transversal estadual, reporta-se a correlação parcial
controlando a latitude}{O gradiente Sul--Norte confunde quase toda variável
observada em Goiás, e sem esse controle a latitude é atribuída à variável que
por acaso a acompanha. A regra nasceu de uma autocorreção, quando uma
associação transversal que parecia identificar um mecanismo próprio, entre a
idade da pastagem e a difusão do plantio direto, encolheu e deixou de ser
sustentável assim que a latitude entrou na mesma conta
(Seção~\ref{sec:res-perna2}).}

\decisao{D15}{O fogo do ano $t$ é alinhado à conversão cuja origem é $t$, e
tratado como contemporâneo}{Alinhá-lo de outro modo atribuiria ao fogo uma
liderança temporal que é em parte definicional, porque a queima faz parte do
próprio evento de conversão que se está medindo.}

\decisao{D16}{Precedência entre séries integradas exige Toda-Yamamoto e
placebos direcionais}{O teste de Granger aplicado a séries integradas fabrica
precedência que não existe, e foi o que aconteceu na primeira versão do teste
de deslocamento (Seção~\ref{sec:met-instrumental}).}

\decisao{D19}{Todo deslocamento de centroide vem com intervalo de confiança de
95\% por \emph{bootstrap}}{As 166 AMCs são reamostradas com reposição, e um
vão entre duas camadas só é lido como real quando as faixas não se sobrepõem.
É o que separa uma diferença medida de duas linhas que parecem distintas no
gráfico.}

\decisao{D20}{Quando o choque é nacional e único, o $p$ vem de permutação, e
não do erro-padrão}{Um choque que é um número por ano para o estado inteiro dá
cerca de quarenta observações independentes, e o erro-padrão agrupado se
comporta como se houvesse milhares. Embaralhar o próprio choque entre os anos
devolve o $p$ que corresponde ao desenho (Seção~\ref{sec:met-instrumental}).}

\section{Quando a amostra virou censo}

As quatro decisões desta seção vêm do mesmo episódio, detalhado na
Seção~\ref{sec:met-d21}.

\decisao{D21}{Toda amostra espacial declara que fração dela caiu fora do
recorte pretendido}{A amostra original era sorteada sobre o retângulo que
envolve Goiás, e não sobre o polígono do estado, de modo que 43,7\% dos pixels
coletados estavam fora do território. O rótulo de cada ponto era corrigido
depois, mas a alocação da amostra já havia ocorrido.}

\decisao{D22}{O código sentinela de erro nunca compartilha valor com categoria
real}{Uma classe ausente da tabela de tradução caía no valor de preenchimento
e passava a contar como dado válido, sem que nada no processamento sinalizasse
a falha. Uma falha de configuração precisa interromper a execução, e não virar
número.}

\decisao{D23}{Com censo, o $p$-valor deixa de medir evidência}{Quando o número
de observações é a população inteira, qualquer desvio ínfimo do modelo nulo
produz estatística arbitrariamente grande, e o que ela mede é o tamanho da
população. A robustez passa a ser buscada na estabilidade entre recortes.}

\decisao{D24}{Estatística ponderada é verificada por contrato}{Toda conta
sobre células com peso tem de reduzir exatamente ao caso sem peso quando o
peso é um, e isso é testado, e não presumido. O contrato vale por módulo,
porque quem implementa a própria conta ponderada precisa da própria
verificação.}

\section{A mudança de rótulo do MapBiomas}

As duas decisões desta seção vêm do mesmo episódio, detalhado na
Seção~\ref{sec:met-d25}.

\decisao{D25}{O objeto medido não é constante ao longo da série}{No trecho
final, parte da terra que sai da pastagem passa a ser rotulada como ``Mosaico
de Usos'' em vez de agricultura, enquanto a estatística de área plantada, que
não passa pelo classificador, registra a lavoura crescendo no mesmo período.
Uma transição pode desaparecer da medida sem que nada tenha mudado no campo.}

\decisao{D26}{Medir em intervalo, com uma âncora vinda de fora do
classificador}{Todo número que tenha agricultura no destino é reportado entre
duas réguas, a classe isolada como piso e sua união com o Mosaico como teto, e
comparado a uma fonte que não passa pelo classificador. Considera-se robusta
apenas a conclusão que sobrevive nas duas pontas.}

\section{Auditoria de figuras e de exposições}

\decisao{D27}{Toda figura importada de um script é reauditada antes de
publicar}{O rótulo de uma figura é uma afirmação e envelhece como qualquer
outra, com o agravante de que um arquivo de imagem não é localizável por busca
textual e não aparece na comparação entre versões. A auditoria pergunta se o
título, os eixos, a legenda e o nome de cada série ainda dizem o que a rotina
conclui hoje. O escopo é dado pelas vistas efetivamente publicadas, e não pelo
inventário de arquivos gerados.}

\decisao{D28}{A exposição do \emph{drive} comum é a régua do gradiente, e não
o canal}{A aptidão edafoclimática correlaciona-se $-0{,}44$ com a latitude, e
posta a latitude na mesma regressão ela perde a maior parte da magnitude e
toda a significância. O achado passa a ser enunciado como gradiente
Sul--Norte (Seção~\ref{sec:met-d28}).}

\decisao{D29}{Variável fora do domínio declarado não entra padronizada}{A
variável de esgotamento é uma fração e deveria viver entre zero e um, e no
painel alcançava $-84{,}9$ em 14\% dos pares unidade-ano. Padronizada, passava
a ser dominada por esses poucos valores (Seção~\ref{sec:met-d29}).}
